\documentclass{article}
\title{Boosting}
\author{QwanxingxuA}
\usepackage{ctex}
\usepackage{amsmath}
\usepackage{amssymb}

\begin{document}
\maketitle
    \section{简介}
        Boosting是机器学习中性能很好的方法。从方法应用的角度来看,它解决了两个核心问题：将模型注意力放到分类错误样本上；组合多个模型，并且
        不是平均个组件的功效，而是通过“学习”这一强大的手段组合各个组件。

        Boosting是对决策树的强大改良，在决策树那一章中笔者提到与神经网络的对比，发现决策树一个难点就是无法像神经网络那样反向传播。
        但是Boosting（提升树）是可以类似的，只是提升树采用“梯度上升”的方式。那么为什么决策树和提升树会出现不同呢？笔者将在
        后续解释。

        本文着重探讨Boosting的机器学习逻辑，以数学辅以证明，关于算法的详细步骤省略，建议初次接触的读者参考李航老师书的算法部分。

        本文主要针对分类模型，会涉及回归模型。下面给出分类问题的前提定义。

        定义样本空间$T=\{(x_1,y_1),(x_2,y_2),\cdots,(x_n,y_n)\}$,输入空间$\mathcal{X}$,输出空间$\mathcal{Y}=\{-1,+1\}$。

    \section{加法模型，AdaBoost}
        \subsection{加法模型}
            加法模型顾名思义就是加在一起的模型,即
            \[f(x)=\sum_{m=1}^{M}\alpha_mb(x;\varTheta_m )\]
            其中$b$是基函数，$\varTheta_m$是基函数的参数。

        \subsection{AdaBoost}
            AdaBoost本质依旧是加法模型，只在策略的选择上使用了指数函数的特例，并且其收敛性具有随机的特性。故而取名
            "AdaBoost"。

            定义模型：
            \[f(x)=\sum_{m=1}^{M}\alpha_mG_m(x)\]

            对于当下特点的二分类任务，更严谨的定义是：
            \[F(x)=sign(f(x))\]
            \[f(x)=\sum_{m=1}^{M}\alpha_mG_m(x)\]

            Boosting算法采用递归提升的方法，有点类似于递归树的提升部分。用数学表达：
            \[f_m(x)=f_{m-1}(x)+\alpha_mG_m(x)\]

            下面考虑策略，AdaBoost使用指数损失函数：
            \[L(y,f(x))=\exp(-yf(x))\]

            于是可以构建AdaBoost的优化问题：
            \[\arg\min_{\alpha_m,G_m}\sum_{i=1}^{n}\exp\left[-y_i(f_{m-1}(x)+\alpha_mG_m(x))\right]\]

            令$w_{mi}=\exp(-y_if_{m-1}f(x))$
            
            \[\arg\min_{\alpha_m,G_m}\sum_{i=1}^{n}w_{mi}\exp(-y_iG_m(x))=\sum_{y_i= G_m(x_i)}w_{mi}e^{-\alpha_m}+\sum_{y_i\neq G_m(x_i)}w_{mi}e^{\alpha_m}\]
            
            下面求解上述问题。对于该分类任务，策略是要减少分类错误的概率。故而得到子问题，即$G_m$的解
            \[G_m^*=\arg\min_{G_m}\sum_{i=1}^{n}w_{mi}I(G_m(x)\neq y_i)\]

            这里容易误解，其实就是原始问题的第二部分的等价优化问题。优化问题转换成等价问题是很常见的手段。

            令$e_m=\sum_{i=1}^{n}w_{mi}I(G_m^*(x)\neq y_i)$,对目标函数求导
            \[\frac{\partial L}{\partial \alpha_m}=-e^{-\alpha_m}(1-e_m)+e^\alpha e_m=0\]

            得到
            \[\alpha_m^*=\frac{1}{2}\log\frac{1-e_m}{e_m}\]

            以上便是AdaBoost的算法过程，其实Adaboost本身就是模型、策略、算法，只是平常我们更多是使用的算法。
            笔者为了迎合这种现象，现将上述求解的结果整理成下面的算法组件。
            \[G_m=\arg\min e_m\]
            \[e_m=\sum_{i=1}^{n}w_{mi}I(G_m(x)\neq y_i)\]
            \[\alpha_m=\frac{1}{2}\log\frac{1-e_m}{e_m}\]
            \[w_{mi}=\exp(-y_if_{m-1}(x_i))=\frac{w_{m-1,i}}{Z_m}\exp(-y_i\alpha_{m-1}G_{m-1}(x_i))\]
            \[Z_m=\sum_{i=1}^{n}\exp(-y_i\alpha_{m-1}G_{m-1}(x_i))\]
            其中$Z_m$是规范化因子。

        \subsection*{AdaBoost训练误差上界}
        AdaBoost训练误差上界为:
        \[\frac{1}{n}\sum_{i=1}^{n}I(F(x_i)\neq y_i)\le \frac{1}{n}\sum_{i=1}^{n}\exp(-y_if(x_i))=\prod_{m=1}^{M}Z_m \]

        证明：
        \begin{enumerate}
            \item 先证明前半部分：
            
            对于分类错误样本有$-y_if(x_i)\ge 0$,则
            \[e^{-y_if(x_i)}\ge 1\]
            \[\frac{1}{n}\sum_{i=1}^{n}I(F(x_i)\neq y_i)\le \frac{1}{n}\sum_{i=1}^{n}\exp(-y_if(x_i))\]
            证毕
            \item 再证明后半部分：
            \begin{align*}
                \frac{1}{n}\sum_{i=1}^{n}\exp(-y_if(x_i))=&\frac{1}{n}\sum_{i=1}^{n}\exp(-y_i\sum_{m=1}^{M}\alpha_mG_m(x_i))\\
                =&\frac{1}{n}\sum_{i=1}^{n}\prod_{m=1}^{M}\exp(-y_i\alpha_mG_m(x_i))\\
                =&\frac{1}{n}\sum_{i=1}^{n}\frac{w_{M+1}}{w_1}\prod_{m=1}^{M}Z_m
            \end{align*}
            由于$w_1=\frac{1}{n},\sum_{i=1}^{n}w_{M+1}=1$
            \[\frac{1}{n}\sum_{i=1}^{n}\exp(-y_if(x_i))=\prod_{m=1}^{M}Z_m\]
            证毕
        \end{enumerate}

        针对二分类问题，上式还可以扩展：
        \begin{align*}
            Z_m=&\sum_{y_i=G_m(x_i)}w_{mi}e^{-\alpha_m}+\sum_{y_i\neq G_m(x_i)}w_{mi}e^{\alpha_m}\\
            =&(1-e_m)e^{-\alpha_m}+e_me^{\alpha_m}\\
            =&2\sqrt{e_m(1-e_m)}
        \end{align*}
        
        令$\gamma_m=\frac{1}{2}-e_m$，由于$e^{-x}\ge\sqrt{1-2x}$
        \[Z_m=\sqrt{1-4\gamma_m^2}\le e^{-2\gamma_m^2}\]

        所以
        \[\prod_{m=1}^{M}Z_m\le \exp(-2\sum_{m=1}^{M}\gamma_m^2)\]

        如果存在$\gamma>0$且对所有的$m$均有$\gamma_m \ge \gamma$,则
        \[\prod_{m=1}^{M}Z_m\le e^{-2M\gamma^2}\]

        这个表达式说明了算法呈指数收敛，这是说明算法具有非常好的性能。由于$\gamma$具有随机性，所以
        算法取名Adaboost。

    \section{提升树}
        \subsection{提升树}
        提升树是二叉树，本质也是加法模型。与Adaboost在策略上具有差异，但是在平方误差下和CART算法类似。

        假设提升树得到$J$个叶子节点$\{R_1,R_2,\cdots,R_J\}$,定义提升树模型：
        \[f(x)=\sum_{m=1}^{M}T(x;\varTheta_m)\]
        \[f_m(x)=f_{m-1}(x)+T(x;\varTheta_m)\]
        \[T(x;\varTheta_m)=\sum_{j=1}^{J}c_jI(x\in R_j)\]

        这个定义和CART回归树是类似的。
        
        取平方误差损失函数：
        \[L(y,f_m(x))=(y-f_{m-1}(x)-T(x;\varTheta_m))^2=(r-T(x;\varTheta_m))^2\]

        其中，$r$是上一子树的残差。

        提升树是二叉树,定义
        \[R_1=\{x\mid  x\le s\},R_2=\{x\mid x >s\}\]

        则问题转化为
        \[\min_{s}\min_{c_1,c_2}[(r_1-c_1)^2+(r_2-c_2)^2]\]

        这个问题也同CART算法类似，具体不再赘述。

        当取一些其它损失函数的时候很难之间求解问题，但是提升树的好处就在于是可以求导传递的，只不过是以
        梯度提升的方法。这也说明树也可以构建微分网络。

        \subsection{提升树与决策树}
        就上面提升树回归树例子来看，提升树与决策树CART简直一样，但实际上这是不对的。

        CART回归树如笔者在决策树说的，它是单层模型，虽然各层次之间存在联系，但是模型本身是忽略了这种联系的。
        严格的说，决策树基于单层次的极大似然或者说最小条件熵构建，本质依旧只是单层模型。提升树的加法模型是整体的，
        这种层次间的模型是严格出现在模型表达式上面的。就CART回归树来看，它是基于标签与叶子节点的损失，但是提升树
        是残差与叶子节点的损失。这是二者在平方损失误差下本质的区别，残差记录了各层次之间的联系，但是CART只是浅显
        的受使用过的特征联系而已。

        提升树是决策树的强大改良。笔者在决策树那一章提到决策树的神经网络构想，其实现在想来是不成立的，因为决策树的
        网络联系是浅显的，这是模型本身决定的，难以传播也是必然的。提升树是有可能的，因为加法模型本身就具体以上
        特性。

        实际上CART剪枝是比较生硬的，它没有神经网络那样基于交叉熵损失函数反向传播优化的便利。
        同时，它也没有提升树这样基于损失函数逐步提升的“传播”。笔者在决策树那一章说剪枝如图神经网络的softmax with 
        crossentropyloss层，现在看来是完全达不到的。
\end{document}